% =====================================================================
\begin{thebibliography}{99}
\addcontentsline{toc}{chapter}{Bibliografía}
\markboth{Bibliografía}{}
\small

\bibitem{apostol}
T.~M. Apostol.
\newblock \emph{Calculus}, vol.~1: \emph{One-Variable Calculus, with an Introduction to Linear Algebra}, 2.ª ed.
\newblock Wiley, 1967.
\newblock Cap.~3 «Continuous Functions», cap.~4 «Differential Calculus», cap.~6 «The Logarithm, the Exponential, and the Inverse Trigonometric Functions» y cap.~7 «Polynomial Approximations to Functions».

\bibitem{cambA1}
W.~T. Gowers (profesor); D.~Chua (transcripción).
\newblock \emph{Part IA Analysis~I}, apuntes de clase.
\newblock University of Cambridge, Mathematical Tripos, Lent 2015.
\newblock Cap.~2 «Convergence of sequences» (§2.1--2.3), cap.~4 «Continuous functions» (§4.1) y cap.~5 «Differentiability» (§5.1--5.3).
\newblock \url{https://dec41.user.srcf.net/h/IA_L/analysis_i}

\bibitem{tripos2013}
University of Cambridge.
\newblock \emph{Mathematical Tripos Part IA}, Paper~1, 30 de mayo de 2013.
\newblock Preguntas 9D, 10E y 11E (Analysis~I).
\newblock \url{https://www.maths.cam.ac.uk/undergrad/pastpapers/files/2013/PaperIA_1.pdf}

\bibitem{tripos2019}
University of Cambridge.
\newblock \emph{Mathematical Tripos Part IA}, Paper~1, 2019.
\newblock Preguntas 3E, 9D y 10D (Analysis~I).
\newblock \url{https://www.maths.cam.ac.uk/undergrad/pastpapers/files/2019/paperia_1_2019.pdf}

\bibitem{mit1801}
D.~Jerison.
\newblock 18.01SC: Single Variable Calculus.
\newblock Massachusetts Institute of Technology, MIT OpenCourseWare, otoño de 2010. Unidad~1 «Differentiation» (sesiones~1--20 y \emph{Problem Sets}~1--2) y unidad~2 «Applications of Differentiation» (sesiones~23--36 y \emph{Problem Sets}~3--5).
\newblock \url{https://ocw.mit.edu/courses/18-01sc-single-variable-calculus-fall-2010/}

\bibitem{mit1801ex}
D.~Jerison y A.~Mattuck.
\newblock \emph{18.01 Exercises}: secciones 1C--1G y~1J («Differentiation») y 2A, 2B, 2C y~2G («Applications of Differentiation»); exámenes~1 y~2 del curso 18.01SC con sus soluciones.
\newblock MIT OpenCourseWare, 2010.
\newblock \url{https://ocw.mit.edu/courses/18-01sc-single-variable-calculus-fall-2010/}

\bibitem{osc1}
G.~Strang y E.~Herman.
\newblock \emph{Calculus}, vol.~1.
\newblock OpenStax, Rice University, 2016.
\newblock Cap.~1 «Functions and Graphs» (§1.1--1.5), cap.~2 «Limits» (§2.1--2.5), cap.~3 «Derivatives» (§3.1--3.9) y cap.~4 «Applications of Derivatives» (§4.1--4.8).
\newblock \url{https://openstax.org/details/books/calculus-volume-1}

\bibitem{osc2}
G.~Strang y E.~Herman.
\newblock \emph{Calculus}, vol.~2.
\newblock OpenStax, Rice University, 2016.
\newblock §6.3 «Taylor and Maclaurin Series».
\newblock \url{https://openstax.org/details/books/calculus-volume-2}

\bibitem{spivak}
M.~Spivak.
\newblock \emph{Calculus}, 4.ª ed.
\newblock Publish or Perish, 2008.
\newblock Cap.~3--12 (de «Functions» a «Inverse Functions»), cap.~18 «The Logarithm and Exponential Functions», cap.~20 «Approximation by Polynomial Functions» y cap.~22 «Infinite Sequences».

\bibitem{stewart}
J.~Stewart.
\newblock \emph{Calculus: Early Transcendentals}, 8.ª ed.
\newblock Cengage Learning, 2016.
\newblock Cap.~1 «Functions and Models» (§1.1--1.5), cap.~2 «Limits and Derivatives» (§2.1--2.8), cap.~3 «Differentiation Rules» (§3.1--3.11), cap.~4 «Applications of Differentiation» (§4.1--4.5 y~4.7) y cap.~11 «Infinite Sequences and Series».

\end{thebibliography}
